\documentclass[11pt,a4paper]{article}
\usepackage[utf8]{inputenc}
\usepackage[french]{babel}
\usepackage[T1]{fontenc}
\usepackage{lmodern}
\usepackage[margin=1.5cm]{geometry}

\begin{document}

% En-tête
\begin{center}
    {\Huge \textbf{Léo Robert}} \\[0.2cm]
    {\Large \textbf{Programmeur Moteur 3D}} \\[0.1cm]
    léo.robert@email.com \ | \ 06 20 33 91 52
\end{center}

\vspace{0.3cm}

% Résumé / Profil
\section*{Profil Professionnel}
\noindent
Ingénieur spécialisé dans le développement de jeux vidéo et la programmation 3D. Passionné par l'optimisation des performances en temps réel et l'architecture de moteurs de jeu, je possède une excellente maîtrise du C++ et des outils de l'industrie (Unreal Engine). Mon approche mathématique et ma rigueur me permettent de repousser les limites visuelles et techniques des plateformes modernes.

\vspace{0.3cm}

% Compétences
\section*{Compétences Techniques}
\noindent
\textbf{Vulkan} $\bullet$ \textbf{C++} $\bullet$ \textbf{Unreal Engine 5} $\bullet$ \textbf{Mathématiques 3D} $\bullet$ \textbf{Git} $\bullet$ \textbf{OpenGL} $\bullet$ \textbf{Perforce} $\bullet$ \textbf{Unity 3D}

\vspace{0.3cm}

% Expériences
\section*{Expériences Professionnelles}
\noindent \textbf{Programmeur Moteur 3D / Confirmé} --- \textit{Arkane Studios} \hfill \textbf{08/2025 à Aujourd'hui} \\
Développement d'un système de cinématiques temps réel, permettant de synchroniser les animations faciales avec le doublage audio. Intégration d'un moteur physique multi-threadé permettant de simuler des destructions de décors complexes à grande échelle sans perte de framerate. Création d'un système d'intelligence artificielle (Behavior Trees) pour les PNJ, rendant leurs réactions plus réalistes et adaptatives face au joueur. Refactoring complet du système d'inventaire modulaire, le rendant compatible avec la sauvegarde cloud multi-plateformes. \\[0.3cm]
\noindent \textbf{Programmeur Moteur 3D} --- \textit{Arkane Studios} \hfill \textbf{12/2023 à 08/2025} \\
Optimisation de la gestion de la mémoire et réduction du temps de chargement des niveaux de 40\% grâce au streaming asynchrone des assets 3D. Développement d'outils internes pour l'éditeur de niveau, augmentant de 30\% la productivité de l'équipe de Level Design au quotidien. Développement de mécaniques de gameplay en C++ et optimisation du rendu graphique sur Unreal Engine pour un titre AAA, garantissant 60 FPS constants. Intégration d'outils de télémétrie in-game pour récolter les données de parcours des joueurs et ajuster la difficulté dynamiquement. \\[0.3cm]
\noindent \textbf{Stage — Assistant Programmeur Moteur 3D} --- \textit{Quantic Dream} \hfill \textbf{04/2022 à 12/2023} \\
Implémentation de shaders personnalisés et de systèmes de particules avancés avec Niagara pour améliorer la qualité visuelle globale du jeu. Portage complet du moteur de jeu vers les consoles nouvelle génération (PS5, Xbox Series X), en exploitant au maximum les architectures SSD. Profilage CPU/GPU intensif à l'aide de PIX et RenderDoc pour identifier et corriger les goulots d'étranglement avant la phase de certification console. \\[0.3cm]


\vspace{0.3cm}

% Formations
\section*{Formation \& Diplômes}
\noindent \textbf{Licence Professionnelle en Informatique Visuelle} \hfill \textbf{2022 - 2023} \\
\textit{Université de Montpellier} \\[0.3cm]
\noindent \textbf{BTS en Informatique Visuelle} \hfill \textbf{2020 - 2022} \\
\textit{IUT Bobigny} \\[0.3cm]


\end{document}